\section{Introduction}
\label{sec:introduction}

On a smooth Riemannian manifold, an upper bound for the Hessian of a
function bounds its second derivative along every geodesic. The same
question is subtler when the function has only a distributional
Hessian. A weak inequality is tested against the reference measure,
whereas a prescribed geodesic can lie in a set of measure zero. The
goal of this paper is to pass from such a weak inequality to a bound
on every minimizing geodesic in a finite-dimensional \(\RCD\) space.

We work on a complete separable metric measure space \((M,d,m)\)
with full support and measure finite on bounded sets. We use
\(\RCD(K,N)\) to mean the unreduced curvature-dimension condition
\(\CD(K,N)\), together with quadratic Cheeger energy. The parameters
satisfy \(K\in\mathbb R\) and \(1<N<\infty\).
The Laplacian is the nonpositive \(L^2\) generator:
\[
\int_M\langle\nabla u,\nabla v\rangle\dd m
=-\int_M(\Delta u)v\dd m.
\]
Its domain \(D(\Delta)\) is contained in \(W^{1,2}(M)\).
Write \(\Gamma(u,v)=\langle\nabla u,\nabla v\rangle\) and
\(\Gamma(u)=\Gamma(u,u)\). Our global test class is
\begin{equation}
\label{eq:test-class}
\Test(M)=\left\{g\in D(\Delta):
\begin{array}{l}
 g\text{ has a bounded globally Lipschitz representative},\\
 \Delta g\in W^{1,2}(M)
\end{array}\right\}.
\end{equation}

The second-order \(\RCD\) calculus assigns an \(L^2\) Hessian to a
test function and gives
\(\nabla_{\nabla g}\nabla g=\nabla(\Gamma(g)/2)\)
\cite{Gigli2018}. For \(F\in W^{1,2}_{\rm loc}(M)\), a compactly
supported \(g\in\Test(M)\), \(Y=\nabla g\), and
\(h\in\Lip_c(M)\), define its weak Hessian evaluation by
\begin{equation}
\label{eq:weak-hessian}
H_F(Y,Y)(h)
=-\int_M\divg(hY)\,\langle\nabla F,Y\rangle\dd m
-\int_M h\,\langle\nabla F,\nabla_YY\rangle\dd m.
\end{equation}
Here \(\divg(hY)=\Gamma(h,g)+h\Delta g\), and both integrals
are finite: the relevant derivatives are locally \(L^2\), and
\(h,Y\) are bounded with compact support. This agrees with the
usual \(L^2\) Hessian when \(F\in D(\Delta)\). The definition
itself requires only first derivatives of \(F\).

\begin{theorem}[Weak Hessian bounds along every geodesic]
\label{weight:every-geodesic}
Let \((M,d,m)\) be a full-support \(\RCD(K,N)\) space,
where \(K\in\mathbb R\) and \(1<N<\infty\). Let
\(F\colon M\to\mathbb R\) be bounded and globally Lipschitz,
and let \(G\colon M\to\mathbb R\) be bounded and continuous.
Suppose that for every compactly supported \(g\in\Test(M)\)
and every nonnegative \(h\in\Lip_c(M)\),
\begin{equation}
\label{weight:hessian-assumption}
H_F(\nabla g,\nabla g)(h)
\le\int_M hG\Gamma(g)\dd m.
\end{equation}
Then every constant-speed minimizing geodesic
\(\sigma\colon[0,1]\to M\), of length \(\ell\), satisfies
\begin{equation}
\label{weight:curve-hessian}
(F\circ\sigma)''\le\ell^2G\circ\sigma
\qquad\text{in distributions on }(0,1).
\end{equation}
\end{theorem}

Thus, for every nonnegative \(\varphi\in C_c^\infty(0,1)\), the
conclusion means
\[
\int_0^1 F(\sigma_t)\varphi''(t)\dd t
\le \ell^2\int_0^1 G(\sigma_t)\varphi(t)\dd t.
\]
For constant \(G=c\), it says that
\(t\mapsto F(\sigma_t)-c\ell^2t^2/2\) is concave.
The ambient reference measure can have infinite total mass:
bounded Lipschitz \(F\) is locally Sobolev, and no global
\(L^2(m)\) condition on \(F\) is needed.

\subsection{Background and the analytic obstacle}

The differential calculus of \(\RCD\) spaces makes weak gradients,
Hessians and integration by parts available without a smooth atlas
\cite{Gigli2018}. Savar\'e's self-improvement of the
Bakry--\'Emery condition supplies a measure-valued Bochner calculus
and the multiplier properties used below \cite{Savare2014}.
Ambrosio--Gigli--Savar\'e identify the corresponding constant
Bakry--\'Emery bounds with Riemannian curvature-dimension bounds
\cite{AGS2015}; Braun--Habermann--Sturm develop the analogous
variable-curvature equivalences \cite{BHS2021}.

Changes of measure and large-parameter entropy limits already play
an important role in Hessian-to-convexity arguments. Ketterer uses
an exponential weight and a scaling limit under a Hessian assumption
that includes Sobolev regularity of gradient pairings
\cite[Theorem~7.1 and its proof]{Ketterer2015Obata}. Sturm develops a
vanishing-entropy argument for semiconvex functions
\cite[Lemmas~1 and 2]{Sturm2018}. Particularly close to our transport
step is Han's proof of infinitesimal-to-weak convexity: he applies
entropy convexity for \(e^{-a u}m\), divides by \(a\), and lets
\(a\) tend to infinity \cite[Theorem~3.13, Equation~(3.6)]{Han2018Monotonicity}.
His hypotheses include local second-order Sobolev regularity of
\(u\). The work below establishes the measure-change step under
the distributional hypothesis \eqref{weight:hessian-assumption}
and retains the spatially varying bound \(G\).

Brena--Gigli study fine measure representations of weak Hessians
and discuss the passage from a weak Hessian bound to geodesic
convexity \cite[Remarks~2.3 and 2.4]{BrenaGigli2024}.
In particular, their Remark~2.4 proposes using a change of measure
and a weak Bochner inequality, and identifies the regularity needed
to justify the approximation in that argument. This is a close
antecedent of the method here. Under our bounded Lipschitz
hypotheses, we prove the multiplier regularity needed for the
measure-change argument by a resolvent contraction for the original
heat semigroup.

For bounded globally Lipschitz \(V\), the change of measure
\(m_V=e^V m\) preserves the \(L^2\) generator domain,
but it need not preserve the test class. Its generator has the drift term
\(\Delta_Vu=\Delta u+\Gamma(V,u)\); a Lipschitz \(V\) does not
ensure that \(\Gamma(V,u)\) is Sobolev. The proof therefore closes
the weighted Bochner inequality in the common generator graph
domain before applying any curvature-dimension equivalence.
The example in Remark~\ref{weight:domain-warning} makes this
regularity issue explicit.

With Laplacian-domain regularity, Kapovitch--Ketterer also pass
from an \(L^2\)-Hessian bound to a geodesic convexity inequality
by approximating point endpoints with measures
\cite[Theorem~4.7]{KapovitchKetterer2020}.
The final passage to a prescribed geodesic uses Deng's
nonbranching theorem for finite-dimensional \(\RCD(K,N)\) spaces
\cite[Theorem~1.3]{Deng2020}. It makes every strict interior
subsegment of that geodesic the unique minimizer between its
endpoints. This metric fact allows entropy inequalities for
concentrating measures to recover the chosen subsegment, including
subsegments invisible to an almost-everywhere assertion about paths.

\subsection{Proof structure}

The proof has two stages. First fix \(\lambda>0\) and change the
measure to \(m_\lambda=e^{\lambda F}m\). The weak Hessian bound
implies a weighted Bochner inequality with curvature
\(K-\lambda G\). A resolvent estimate gives Lipschitz regularity
of the weighted multipliers, and the common graph core makes the
Bochner inequality valid on the actual weighted generator domain.
The constant and variable curvature-dimension equivalences then
yield \(\CD(K-\lambda G,\infty)\).

Second, fix a strict subsegment of the geodesic and approximate its
endpoints by normalized measures on small balls. All transport
geodesics between these shrinking supports converge to that unique
subsegment. The identity for Boltzmann entropy
\[
\Ent_{m_\lambda}(\mu)
=\Ent_m(\mu)-\lambda\int F\dd\mu
\]
then turns the variable-curvature inequality into a lower bound for
\(F\) above its endpoint chord, with the Green-kernel correction
determined by \(G\). Choosing \(\lambda\) to dominate the endpoint
entropies removes the remaining entropy term. Applying this bound
on every strict subinterval proves the distributional inequality.

Section~\ref{sec:preliminaries} gives the common calculus and
approximation facts. Section~\ref{sec:drift} establishes the weighted
generator and resolvent estimates. Section~\ref{sec:curvature-change}
proves the variable curvature-dimension bound, and
Section~\ref{sec:every-geodesic} completes the localization argument.
